% =============================================================
%  AI-Powered Robo-Advisor Platform
%  Academic Report — Programming in Finance II (2026)
%  USI — Prof. Peter H. Gruber
%  GitHub: https://github.com/<YOUR_ORG>/robo-advisor
% =============================================================
\documentclass[11pt, a4paper]{article}

% ── Packages ─────────────────────────────────────────────────
\usepackage[T1]{fontenc}
\usepackage[utf8]{inputenc}
\usepackage[english]{babel}
\usepackage{lmodern}
\usepackage{microtype}

\usepackage{geometry}
\geometry{margin=2.5cm}

\usepackage{amsmath, amssymb}
\usepackage{booktabs}       % professional tables
\usepackage{graphicx}
\usepackage{xcolor}
\usepackage{listings}       % code blocks
\usepackage{hyperref}
\usepackage{cleveref}
\usepackage{float}
\usepackage{caption}
\usepackage{subcaption}
\usepackage{enumitem}
\usepackage{csquotes}
\usepackage[style=authoryear, backend=biber]{biblatex}
\addbibresource{references.bib}

% ── Code listing style ────────────────────────────────────────
\lstset{
  basicstyle=\ttfamily\small,
  keywordstyle=\color{blue!70},
  commentstyle=\color{gray},
  stringstyle=\color{teal},
  breaklines=true,
  frame=single,
  numbers=left,
  numberstyle=\tiny\color{gray}
}

% ── Metadata ─────────────────────────────────────────────────
\title{%
  \textbf{AI-Powered Robo-Advisor Platform}\\[0.4em]
  \large Academic Project Report\\
  Programming in Finance II — USI, 2026
}
\author{%
  P1 (Backend/Data) \quad P2 (Quant/Optimizer)\\
  P3 (ML/Profiling) \quad P4 (Frontend/LLM/Docs)
}
\date{May 2026}

% =============================================================
\begin{document}
% =============================================================

\maketitle
\tableofcontents
\newpage

% ─────────────────────────────────────────────────────────────
\section{Introduction}
\label{sec:intro}
% Owner: P4 (integra contributi del team)
% TODO W4: problem statement, gap in existing robo-advisors,
%          our contribution, GitHub link
% ─────────────────────────────────────────────────────────────

\subsection{Problem Statement}

Retail investors increasingly rely on algorithm-driven portfolio recommendations, yet most commercial robo-advisors offer limited transparency regarding their underlying models, data sources, and geographic limitations. This project addresses the gap between quantitative rigour and user-facing explainability.

\subsection{Our Contribution}

% TODO: fill after demo is stable — 3–4 sentences on what
%       the platform does and what makes it distinctive

We present an educational robo-advisor prototype that combines:
\begin{itemize}
  \item Hierarchical Risk Parity (HRP) portfolio optimisation with
        Ledoit--Wolf shrinkage on a UCITS-aware ETF universe;
  \item A Gradient Boosting (GBM) risk profiler trained on real
        household financial behaviour (Fed SCF~2022) with SHAP
        explainability;
  \item A constrained LLM narrator layer (Claude API) that translates
        numerical results into natural language without inventing figures;
  \item An EU-awareness layer that explicitly communicates the US-centric
        bias of the training data to European investors.
\end{itemize}

The source code is available at:
\url{https://github.com/<YOUR_ORG>/robo-advisor}

% ─────────────────────────────────────────────────────────────
\section{ML Risk Profiler}
\label{sec:ml}
% Owner: P3 (draft) — P4 integrates
% TODO W4: GBM on SCF 2022, SHAP, clustering, feature importance
% ─────────────────────────────────────────────────────────────

\subsection{Data: Federal Reserve SCF 2022}
% TODO P3: dataset description, variables used, preprocessing

\subsection{Label Construction via Behavioural Clustering}
% TODO P3: k-means / hierarchical clustering on financial behaviour,
%          why not Grable-Lytton (circular), silhouette score

\subsection{GBM Classifier}
% TODO P3: model architecture, hyperparameters, train/test split,
%          accuracy / F1

\subsection{SHAP Explainability}
% TODO P3: top drivers, waterfall plots, how SHAP values feed into
%          the Ground Truth JSON

% ─────────────────────────────────────────────────────────────
\section{Portfolio Optimisation}
\label{sec:portfolio}
% ─────────────────────────────────────────────────────────────

\subsection{ETF Universe}

The investment universe comprises eight ETFs organised into four
structural clusters reflecting their expected correlation behaviour
under both normal and stress market regimes (Table~\ref{tab:universe}).

\begin{table}[h]
\centering
\caption{ETF Universe v3.1 — UCITS-aware hybrid configuration}
\label{tab:universe}
\small
\begin{tabular}{llllc}
\toprule
\textbf{Primary Ticker} & \textbf{Fallback} & \textbf{Asset Class} & \textbf{Cluster} & \textbf{UCITS} \\
\midrule
CSPX.L & SPY & Equity USA (S\&P 500)         & Risk Assets  & \checkmark \\
EFA    & EFA & Equity International (EAFE)   & Risk Assets  &            \\
GLD    & GLD & Gold                          & Real Assets  &            \\
VNQ    & VNQ & Real Estate (US REITs)        & Real Assets  &            \\
AGGH.MI & AGG & Bond Aggregate (EUR-hedged)  & Safe Haven   & \checkmark \\
TLT    & TLT & US Treasury 20Y+              & Safe Haven   &            \\
TIP    & TIP & US TIPS (Inflation-Linked)    & Safe Haven   &            \\
XEON.MI & BIL & Cash / EUR Overnight Rate   & Cash         & \checkmark \\
\bottomrule
\end{tabular}
\end{table}

Three of the eight ETFs use UCITS-compliant primary tickers
(CSPX.L, AGGH.MI, XEON.MI), improving eligibility for EU retail
investors under MiFID~II. For each ticker, a US-listed fallback is
defined: the \texttt{ValidatedDataLoader} activates the fallback
automatically when the primary ticker exhibits a NaN rate above 2\%,
logging the substitution in the audit trail.

\subsection{Hierarchical Risk Parity}

Portfolio allocation follows the Hierarchical Risk Parity (HRP)
algorithm of \textcite{LopezdePrado2016}, which constructs a
diversified portfolio without requiring an invertible covariance
matrix or explicit return estimates. The algorithm proceeds in three
steps.

\paragraph{Step 1 — Hierarchical Clustering.}
A distance matrix is derived from the shrunk correlation matrix:
\[
  D_{ij} = \sqrt{\tfrac{1}{2}(1 - \rho_{ij})}
\]
Ward linkage is applied to the condensed distance matrix, producing a
dendrogram that reflects the hierarchical correlation structure of the
assets. The resulting tree is visualised in the dashboard as an
interactive Plotly dendrogram.

\paragraph{Step 2 — Quasi-Diagonalisation.}
Assets are reordered according to the leaf sequence of the dendrogram,
placing correlated assets adjacent to one another. This reordering
renders the covariance matrix approximately block-diagonal, exposing
the cluster structure used in the subsequent allocation step.

\paragraph{Step 3 — Recursive Bisection.}
Capital is allocated by recursively splitting each cluster into two
sub-clusters and assigning weights inversely proportional to the
intra-cluster variance:
\[
  \alpha_k = \frac{\sigma^2_{R}}{\sigma^2_{L} + \sigma^2_{R}}, \qquad
  \text{where } \sigma^2_C = \mathbf{w}^{\text{IVP}\top}_C \,
  \Sigma_C \, \mathbf{w}^{\text{IVP}}_C
\]
and $\mathbf{w}^{\text{IVP}}$ denotes the inverse-variance portfolio
within the cluster. This procedure does not require matrix inversion
and is therefore numerically stable on small asset universes.

\subsection{Covariance Estimation: Ledoit--Wolf Shrinkage}

Sample covariance matrices estimated on short return histories are
notoriously noisy: eigenvalues are systematically biased, leading to
extreme and unstable portfolio weights \parencite{Michaud1989}. We
address this by applying the analytical Ledoit--Wolf shrinkage
estimator \parencite{LedoitWolf2004}:

\[
  \hat{\Sigma} = (1 - \alpha)\,S + \alpha\,\mu_S\,I
\]

where $S$ is the sample covariance matrix, $\mu_S$ its average
diagonal entry, $I$ the identity matrix, and $\alpha \in [0,1]$ the
shrinkage intensity estimated analytically (Oracle Approximating
Shrinkage). The estimator pulls extreme eigenvalues toward a common
value, reducing estimation error without discarding the correlation
structure present in the data.

The same shrinkage estimator is applied to both the HRP and the
Mean--Variance benchmark, ensuring that any performance differences
between the two algorithms reflect optimisation methodology, not
input quality.

\subsection{Profile Tilt and Box Constraints}

Raw HRP weights are adjusted in two stages to incorporate investor
risk preferences.

\paragraph{Profile Tilt.}
A 30\% blend is applied between the HRP weights and a
profile-specific target:
\begin{itemize}
  \item \textit{Conservative}: blend toward Minimum Volatility weights.
  \item \textit{Moderate}: HRP weights retained unmodified.
  \item \textit{Aggressive}: blend toward Equal Risk Contribution (ERC) weights.
\end{itemize}

\paragraph{Box Constraints.}
After tilting, weights are clipped to enforce:
\begin{itemize}
  \item Asset level: $[0.05,\; 0.40]$ per instrument.
  \item Cluster level: $[0.10,\; 0.60]$ per structural cluster.
\end{itemize}
Clipping is iterated up to ten times with renormalisation at each
step to ensure convergence. If any constraint is active, the
\texttt{solver\_status} field of \texttt{OptimizationResult} is set
to \texttt{"clipped"}.

\subsection{Mean--Variance Comparison}

The dashboard exposes a dedicated MV tab displaying the efficient
frontier alongside both the HRP and the Max-Sharpe Markowitz
portfolios. The comparison is educational rather than normative:
HRP is the default allocation method; MV serves as an academic
benchmark.

Expected returns for MV are estimated via historical mean log
returns --- a known limitation acknowledged in \textcite{Michaud1989},
who demonstrated that estimation error in $\mu$ causes MV to produce
highly concentrated, out-of-sample unstable portfolios. The same
Ledoit--Wolf covariance is used for both methods. The risk-free rate
is set to 3\% per annum, reflecting approximate EUR short-term rates
in 2024.

The \texttt{compare\_hrp\_vs\_mv()} helper function returns a
structured dictionary consumed by the Streamlit frontend to render a
side-by-side weight comparison table and a weight-divergence bar
chart, making the methodological trade-offs transparent to the user.

% ─────────────────────────────────────────────────────────────
\section{LLM Narrator and Validator}
\label{sec:llm}
% Owner: P4
% TODO W4: narrator pattern, GT JSON, validator 4-step,
%          prompt injection defense, EU Awareness Rule 9
% ─────────────────────────────────────────────────────────────

\subsection{Design Principle: Narrator, Not Calculator}

A language model is not a reliable arithmetic engine. Our architecture
enforces a strict separation: all numerical computations are performed
by the quantitative backend; the LLM receives a \emph{Ground Truth JSON}
containing pre-computed figures and is permitted only to narrate them.

\subsection{Ground Truth JSON Contract}
% TODO W4: schema description, key fields, invariants
% Reference: docs/ground_truth_schema.md

The Ground Truth JSON payload passed to the LLM at inference time
includes the following top-level blocks:

\begin{itemize}
  \item \texttt{portfolio}: weights, risk contributions, expected
        drawdown under stress scenarios;
  \item \texttt{profiler}: \texttt{profile\_label},
        \texttt{profile\_confidence}, \texttt{top\_drivers} (SHAP);
  \item \texttt{regulatory\_context}: MiFID~II disclaimer,
        currency risk note, UCITS eligibility flags,
        SCF US-centrism caveat;
  \item \texttt{llm\_constraints}: \texttt{forbidden\_phrases},
        \texttt{allowed\_numbers}, \texttt{disclaimer\_required}.
\end{itemize}

\subsection{System Prompt Design}

% TODO W4: paste key rules (not full prompt), highlight
%          forbidden phrases, allowed_numbers binding, Rule 9

Key constraints embedded in the system prompt:
\begin{enumerate}
  \item The narrator may only cite numbers present in
        \texttt{allowed\_numbers}.
  \item Phrases such as ``you should invest'', ``guaranteed return'',
        and ``this is financial advice'' are forbidden.
  \item Every response must end with the MiFID~II disclaimer.
  \item \textbf{Rule~9 — EU Awareness}: if the investor is located in
        the EU, the narrator must acknowledge that the risk profiler
        was trained on US household data (Fed SCF) and that results
        should be interpreted with this limitation in mind.
\end{enumerate}

\subsection{Validator: 4-Step Pipeline}

% TODO W4: describe each step, fallback behaviour, test coverage

\begin{enumerate}
  \item \textbf{Number check}: all figures in the response must appear
        verbatim in \texttt{allowed\_numbers}.
  \item \textbf{Forbidden phrase scan}: regex scan against the
        blacklist.
  \item \textbf{Disclaimer presence}: assert disclaimer substring
        present.
  \item \textbf{Length / coherence guard}: response must be between
        50 and 800 tokens; if any step fails, a safe fallback
        message is returned instead.
\end{enumerate}

\subsection{Prompt Injection Defence}

% TODO W4: input sanitisation, max length, character whitelist

% ─────────────────────────────────────────────────────────────
\section{Backtest Results}
\label{sec:backtest}
% Owner: P2 (tables) — P4 integrates
% TODO W4: 3 historical scenarios, Sharpe, max drawdown, table
% ─────────────────────────────────────────────────────────────

\subsection{Scenarios}
% TODO P2: COVID crash (2020), rate hike cycle (2022),
%          2008 stress test if data available

\subsection{Performance Metrics}

% TODO P2: fill table with real numbers after backtest runs

\begin{table}[H]
\centering
\caption{Backtest summary — HRP vs.\ Mean--Variance}
\label{tab:backtest}
\begin{tabular}{lccccc}
\toprule
Scenario & Model & Ann.\ Return & Volatility & Sharpe & Max DD \\
\midrule
COVID-19 (2020) & HRP & \textit{TBD} & \textit{TBD} & \textit{TBD} & \textit{TBD} \\
                & MV  & \textit{TBD} & \textit{TBD} & \textit{TBD} & \textit{TBD} \\
Rate hike (2022)& HRP & \textit{TBD} & \textit{TBD} & \textit{TBD} & \textit{TBD} \\
                & MV  & \textit{TBD} & \textit{TBD} & \textit{TBD} & \textit{TBD} \\
\bottomrule
\end{tabular}
\end{table}

% ─────────────────────────────────────────────────────────────
\section{Limitations and Failure Modes}
\label{sec:limitations}
% Owner: P4 (integra contributi) — all members contribute
% ─────────────────────────────────────────────────────────────

\subsection{Data Fragility: yfinance}
% TODO: discuss single-source dependency, fallback logic

\subsection{SCF US-Centrism vs.\ EU Investor}

The risk profiler is trained on the Federal Reserve Survey of Consumer
Finances~(2022), which captures US household financial behaviour.
European investors may exhibit materially different risk preferences,
savings rates, and financial literacy patterns. The EU Awareness layer
communicates this limitation explicitly at the point of recommendation.

\subsection{HRP Opacity vs.\ Mean--Variance Interpretability}
% TODO: discuss black-box nature of HRP for a retail investor

\subsection{LLM Hallucination Risk}
% TODO: describe residual risk after validator, honest about limits

% ─────────────────────────────────────────────────────────────
\section{Lessons Learned and AI Tools}
\label{sec:lessons}
% Owner: all — P4 coordinates
% TODO W4 Fri–Sun: retrospective, agentic workflow, what worked
% ─────────────────────────────────────────────────────────────

\subsection{Agentic Development Workflow}

% TODO: describe AGENTS.md, GitHub Actions + Claude API for
%       automated PR, evidence of agent PR on GitHub

\subsection{AI Tools Used}

\begin{description}
  \item[Claude (Anthropic)] Used as technical advisor throughout the
    project; generated and reviewed code for
    \texttt{narrator.py}, \texttt{validator.py}, and documentation
    scaffolding.
  \item[GitHub Copilot / ChatGPT] % TODO: add if used
  \item[GitHub Actions] Automated CI (ruff linter, pytest) and
    agent-driven PR generation.
\end{description}

\subsection{What Worked}
% TODO: honest retrospective — 3–5 points

\subsection{What Did Not Work}
% TODO: honest retrospective — 2–3 points

% ─────────────────────────────────────────────────────────────
\section{Conclusions}
\label{sec:conclusions}
% Owner: P4
% TODO W4: 1 page summary, GitHub link, future work
% ─────────────────────────────────────────────────────────────

% TODO: fill after demo is stable

\subsection{Future Work}

\begin{itemize}
  \item Full EU-specific dataset (HFCS) to replace SCF as training
        source for European investors.
  \item Real-time market data integration beyond yfinance.
  \item Regulatory compliance layer for MiFID~II suitability
        assessment (beyond educational disclaimer).
\end{itemize}

% ─────────────────────────────────────────────────────────────
\printbibliography
% ─────────────────────────────────────────────────────────────

\end{document}
